\scanpage{17}
De plus, la construction donne un hom can
\[
E_n\xrightarrow{u_n}E_{n+1}
\]
tel que l'hom correspondant
\diagram{\begin{tikzcd}[ampersand replacement=\&,row sep=0.3em]
u_n\otimes_A B: \& E_n/fE_n\arrow[r,dash]\arrow[d,dash] \& E_{n+1}/fE_{n+1}\arrow[d,dash]\\
\& Q_n \& Q_{n+1}
\end{tikzcd}}
soit \uline{nul} sur $g(E_n/fE_n)$, et induise un hom injectif de \inserted{\unsure{module libre|un mod libre}} $\deleted{E_n/fE_n}Q_n\otimes_B B/gB$ dans \deleted{\unsure{et définissant une base libre|est définie et un hom injec} $Q_{n+1}$ \unsure{en appliquant une base de}} et de $Q_{n+1}\otimes B/gB$, appliquant une base sur une base ce sont des modules libres).

Donc
\[
\begin{aligned}
\varinjlim_n Q_n&=\varinjlim_n Q_n\otimes_B B/gB\simeq\varinjlim (B/gB)^{n-1}\\
&\simeq (B/gB)^{(\mathbb N)}.
\end{aligned}
\]
Prenons donc
\[
M=\varinjlim_n E_n
\]
Alors $M$ satisfait aux conditions voulues.
% archive sheet number: 16

